\documentclass[10pt,a4paper]{amsart}
\usepackage[margin=19mm]{geometry}
\usepackage{amsmath,amssymb,mathtools}
\usepackage[scr=rsfs,cal=euler]{mathalfa}
\usepackage{xfrac}
\usepackage[unicode,hidelinks,bookmarks=false]{hyperref}
\newcommand{\ie}{i.e.\ }
\newcommand{\cf}{cf.\ }
\newcommand{\resp}{respectively\ }
\newcommand{\Subsection}{\subsection}
\providecommand{\nobreakdash}{\nobreak\hbox{-}}
\setlength{\emergencystretch}{2em}
\multlinegap0pt
\usepackage{bm}
\usepackage{bbm}
\usepackage{dsfont}
\usepackage{xspace}
\usepackage{cancel}
\usepackage{mathtools}
\usepackage{amsthm}
\makeatletter
\@ifundefined{theorem}{\newtheorem{theorem}{Theorem}}{}
\@ifundefined{lemma}{\newtheorem{lemma}[theorem]{Lemma}}{}
\makeatother
\title[Existence and multiplicity of normalized solutions for the g]{Existence and multiplicity of normalized solutions for the generalized Kadomtsev-Petviashvili equation in $\mathbb{R}^2$}
\author{Claudianor O. Alves, Rui Ding, Chao Ji}
\date{Source: arXiv:2506.04967v1 (CC BY 4.0)}
\begin{document}
\maketitle

\noindent\emph{Source and licence.} Existence and multiplicity of normalized solutions for the generalized Kadomtsev-Petviashvili equation in $\mathbb{R}^2$. Version arXiv:2506.04967v1 (CC BY 4.0), distributed under the Creative Commons Attribution 4.0 International licence, as recorded in the arXiv licence metadata for this exact version and shown on the abstract page https://arxiv.org/abs/2506.04967v1. Full rights notice: licenses/arxiv-2506.04967-CC-BY-4.0.txt.\par\smallskip

\noindent\textit{Editorial scope.} Counted unit: the Lemma whose source label is \texttt{Lem4-m(a)cont} in the article (source0.tex lines 1224--1226, source label \texttt{Lem4-m(a)cont}), delivered together with the article's own complete original proof of it (source0.tex lines 1227--1294, one \texttt{proof} environment of 68 lines). No mathematics is written, repaired or re-derived: statement and proof are the article's own text line for line. Required context retained uncounted: Eq-Phi>f (lines 1007--1009); lem5-f (lines 1089--1094); Lem5-m(a) (lines 1122--1127). Every definition, display and label of those lines is retained unchanged. Omitted from this package and not redistributed: 1--1006, 1010--1088, 1095--1121, 1128--1223, 1295--1638. Nothing inside a retained excerpt is omitted or altered: every retained body is delivered exactly as the article's lines are written, so no omission receipt of any kind accompanies this package. Citations: the retained excerpts carry no citation command at all, so no bibliography entry of the article is transcribed and no reference is redistributed. Reference adaptation: none was needed and none was made; every reference inside the delivered unit resolves inside it, so no unresolved reference or citation is produced. Printed slips kept literal: no printed word, symbol, hypothesis or number of the counted statement or of its proof was altered. Layout and class: the article's own class is \texttt{amsart} with packages including amsmath, amssymb, amsthm, latexsym, cite, cancel, caption, graphicx, wasysym, overpic, tikz, color, subfigure, hyperref; the wrapper is the lane's pinned \texttt{amsart} editorial wrapper at 10pt on A4, which restates the theorem-like environments the retained text needs and adds the article's own macro definitions that the retained text needs (none), with extra wrapper packages (bm, bbm, dsfont, xspace, cancel, mathtools). The delivered numbering is the unit's own. Assets: no retained span contains a figure environment, a graphics-inclusion command, a PDF-page inclusion, a table or a TikZ picture; no image file is copied into this package, the package declares no asset dependency and no image file is redistributed with it. Licence: version arXiv:2506.04967v1 of this article is distributed under the Creative Commons Attribution 4.0 International licence, as recorded in the arXiv licence metadata for this exact version and shown on the abstract page https://arxiv.org/abs/2506.04967v1. A full rights notice is delivered at licenses/arxiv-2506.04967-CC-BY-4.0.txt. Characters: every retained excerpt is pure ASCII, so the delivered engine, XeLaTeX with its default Latin Modern text font, reports no missing character.
  \begin{equation}\label{Eq-Phi>f}
    J_{\mu}(u) \geq {h\left(a,\|u\|_{0}\right)\| u\|_{0}^2}, \quad \text { for all } u \in \mathcal{D}(a).
  \end{equation}

    \begin{lemma}\label{lem5-f}
    Let $\left(a_1, \rho_1\right) \in(0, \infty) \times(0, \infty)$ be such that $h\left(a_1, \rho_1\right) \geq 0$. Then, for any $a_2 \in\left(0, a_1\right]$, we have
    $$
    h\left(a_2, \rho_2\right) \geq 0 \quad \text {for all} \quad \rho_2 \in\left[\frac{a_2}{a_1} \rho_1, \rho_1\right] .
    $$
    \end{lemma}

\begin{lemma}\label{Lem5-m(a)}
    For any $a \in\left(0, a_0\right)$,
    \begin{equation}\label{eq-m<0}
    m(a)=\inf _{u \in V(a)} J_\mu(u)<0<\inf _{u \in \partial V(a)} J_\mu(u).
    \end{equation}
\end{lemma}	

    \begin{lemma}\label{Lem4-m(a)cont}
    {   The map $a \in (0, a_0) \mapsto m(a)$ is continuous.}
    \end{lemma}

    \begin{proof}
        Let $a \in\left(0, a_0\right)$ be arbitrary and $\left(a_n\right) \subset\left(0, a_0\right)$ with $a_n \rightarrow a$ as $n\rightarrow\infty$.
        It suffices to show that
        $m\left(a_n\right) \rightarrow m(a)$.
        {From the definition of $m\left(a_n\right)$, $m\left(a_n\right)<0$ and Lemma \ref{Lem5-m(a)},} for any $\varepsilon>0$ sufficiently small, there exists $u_n \in V\left(a_n\right)$ such that
        \begin{equation}\label{Eq-Phiun<0}
            J_\mu\left(u_n\right) \leq m\left(a_n\right)+\varepsilon \quad \text { and } \quad J_\mu\left(u_n\right)<0.
        \end{equation}
        {We first show that $(u_n)$ is bounded in $X$. Since $\left\| u_n\right\|_{0}<\rho_0$, the sequences {$(\left|u_n\right|_p)$ and $(\left|u_n\right|_{q})$  are bounded  in $\mathbb{R}$,} and from \eqref{Eq-Phiun<0},  $(u_n)$ is also bounded in $X$. Define $v_n:={\frac{a}{a_n}} u_n$, so $v_n \in {S}(a)$. Next, we are going to prove that}
        \begin{equation}\label{Eq-mPhiv}
            m(a) \leq J_\mu\left(v_n\right).
        \end{equation}
        First, we consider case that $a_n \geq a$, then
        \begin{equation}
            \left\| v_n\right\|_{0}=\frac{a}{a_n}\left\| u_n\right\|_{0} \leq\left\| u_n\right\|_{0}<\rho_0.
        \end{equation}
        Since we have that $v_n \in V(a)$, \eqref{Eq-mPhiv} holds.


        Second, we consider case that $a_n \leq a$, by Lemma \ref{lem5-f},
        $ h\left(a, \rho_0\right) > 0 $.
        Due to the continuity of $\rho \rightarrow h(a, \cdot)$, we may assume there exists sufficiently small
        $\delta>0$ such that
        $$
        h\left(a, \rho\right) > 0\quad \text{for}\quad \rho \in\left[\rho_0, (1+\delta)\rho_0\right].
        $$
        Hence, we deduce from \eqref{Eq-Phi>f} and \eqref{Eq-Phiun<0} that
        $\left\| u_n\right\|_0<\frac{a_n}{a}  \rho_0$ and
        for sufficiently large $n$
        $$
        \left\| v_n\right\|_{0}=\frac{a}{a_n}\left\| u_n\right\|_{0}< \rho_0\leq(1+\delta)\rho_0.
        $$
        If $\rho_0\leq\left\|v_n\right\|_{0}\leq(1+\delta)\rho_0$, from
        \eqref{Eq-Phi>f}
        $$
        m(a_n)<0<h\left(c,\|v_n\|_{0}\right)\| v_n\|_{0}^{2}\leq J_\mu\left(v_n\right).
        $$
        If $\left\|v_n\right\|_{0}<\rho_0$, we have that $v_n \in V(a)$, so \eqref{Eq-mPhiv} holds. Based on the above two cases, we can conclude that
        \begin{equation}
            m(a) \leq J_\mu\left(v_n\right)=J_\mu\left(u_n\right)+\left(J_\mu\left(v_n\right)-J_\mu\left(u_n\right)\right)
        \end{equation}
        and
        \begin{equation}
            \begin{aligned}
                J_\mu\left(v_n\right) - J_\mu\left(u_n\right) = & -\frac{1}{2} \left( \frac{a}{a_n} - 1 \right) \left\| u_n\right\|_0^2
                -
                 + \frac{\mu}{q} \left[ \left( \frac{a}{a_n} \right)^{\frac{q}{2}} - 1 \right] \left|u_n\right|_q^q \\
                &+ \frac{1}{p} \left[ \left( \frac{a}{a_n} \right)^{\frac{p}{2}} - 1 \right] \left|u_n\right|_{p}^{p}.
            \end{aligned}
        \end{equation}
        Since $ (u_n) $ is bounded in $X$, {we infer that}
        \begin{equation}\label{Eq-m(a)<Phivn}
            m(a) \leq J_\mu\left(v_n\right)=J_\mu\left(u_n\right)+o_n(1).
        \end{equation}
        Combining \eqref{Eq-Phiun<0} and \eqref{Eq-m(a)<Phivn}, {we arrive at}
        $$
        m(a) \leq m\left(a_n\right)+\varepsilon+o_n(1) .
        $$
        Now, let $u \in V(a)$ be such that
        $$
        J_\mu(u) \leq m(a)+\varepsilon \quad \text { and } \quad J_\mu(u)<0 .
        $$
        Define $w_n:={\frac{a_n}{a}} u$, so $w_n \in S\left(a\right)$. Clearly, $\| u\|_{0}<\rho_0$ and $c_n \rightarrow c$ imply $\left\| w_n\right\|_{0}<\rho_0$ for $n$ large enough, thus $w_n \in V\left(a_n\right)$. Moreover, since $J_\mu\left(w_n\right) \rightarrow J_\mu(u)$, it follows that
        $$
        m\left(a_n\right) \leq J_\mu\left(w_n\right)=J_\mu(u)+\left(J_\mu\left(w_n\right)-J_\mu(u)\right) \leq m(a)+\varepsilon+o_n(1).
        $$
    Since $\varepsilon>0$ is arbitrary, it follows that $m\left(a_n\right) \rightarrow m(a)$, completing the proof.
    \end{proof}

\end{document}
